In this section we work in the field $\widetilde{\C}$ of complex
Puiseux series, considered as an elementary substructure of
$\mathbb U$. Put $F=\C((t))$ and take the parameter base to be
$\mathbb S=F$. Normalize the valuation by $v(t)=1$, so that the
value group of $\widetilde{\C}$ is identified with $\Q$.
The Newton--Puiseux theorem identifies
$\widetilde{\C}=\bigcup_{N\geq1}\C((t^{1/N}))$ as an algebraic closure
of $F$; see \cite[Corollary 13.15]{Eisenbud}.
All definability and all categories are interpreted in the original
language $\lan{}{RV}{}$ over $\mathbb S$. In particular, each
$a/N\in\Q$ is definable over $\mathbb S$ as the unique solution of
$N\gamma=a\,v(t)$.

\begin{rem}[The finite Kummer actions]\label{puiseux:kummer:action}
There are natural identifications
\[
 \gal(\widetilde{\C}/F)\ \cong\ \widehat{\mu}
 \ \cong\ \aut\bigl(\RV(\widetilde{\C})/\RV(F)\bigr),
 \qquad
 \widehat{\mu}=\varprojlim_N\mu_N,
\]
where the transition map $\mu_M\to\mu_N$ for $N\mid M$
is $\zeta\mapsto\zeta^{M/N}$, and the last automorphism group
uses the induced original-language structure on $\RV$.

For the proof, fix $N\geq1$, choose any root $\tau_N$ of
$X^N-t$ in $\widetilde{\C}$, and put
$F_N=\C((\tau_N))$. Grouping exponents modulo $N$ writes each
element of $F_N$ uniquely as
\[
 \sum_{j=0}^{N-1}\tau_N^j f_j(t),\qquad f_j(t)\in F.
\]
Uniqueness also follows from the pairwise distinct valuations
modulo $\Z$ of the nonzero summands. Thus $[F_N:F]=N$.
Its $F$-automorphisms are precisely
$\tau_N\mapsto\zeta\tau_N$ for $\zeta\in\mu_N$.
Replacing $\tau_N$ by another root multiplies it by an element
of $\mu_N\subseteq F$, so both $F_N$ and the ratio
$\sigma(\tau_N)/\tau_N$ are independent of that choice.
For $N\mid M$ these ratios satisfy
\[
 \frac{\sigma(\tau_N)}{\tau_N}
 =
 \left(\frac{\sigma(\tau_M)}{\tau_M}\right)^{M/N},
\]
since $\tau_M^{M/N}/\tau_N\in\mu_N\subseteq F$.
The fields $F_N$ have union $\widetilde{\C}$.
Their finite automorphisms therefore give exactly the compatible
systems $(\zeta_N)_N\in\widehat{\mu}$, proving the first
identification. Each such automorphism preserves the valuation:
in the displayed expansion it multiplies the $j$th summand by
$\zeta_N^j$, preserving its valuation and the unique least
valuation of a nonzero sum.

Write $r_N=\rv(\tau_N)$ and identify $\C^\times$ with the
value-zero subgroup of $\RV(\widetilde{\C})$.
For any $u$ of value $a/N$ there is a unique $c\in\C^\times$
such that $u=c\,r_N^a$. The induced action is
\[
 \sigma_{(\zeta_N)}(c\,r_N^a)
   =c\,\zeta_N^a r_N^a .
\]
In particular, a nonidentity system moves $r_N$ for some $N$.
The action on $\RV$ is faithful.

Finally, $\RV(F)$ is the subgroup generated by
$\C^\times$ and $\rv(t)$, by taking the leading term of a
Laurent series. An automorphism of the induced RV structure
fixing $\RV(F)$ must send $r_N$ to
$\zeta_Nr_N$ for some $\zeta_N\in\mu_N$: these are precisely
the $N$ roots of $\rv(t)$ in $\RV$.
The same relations as above force compatibility of the
$\zeta_N$. The displayed action then determines the
automorphism on every RV element. It agrees with the unique
valued-field automorphism constructed from that system.
This proves the second identification.
The parameter base for all these definability assertions is
still $\mathbb S=F$.
\end{rem}

\begin{nota}[Equivariant residue realization]\label{bipo:gdv}
Write $\ggk^{\widehat\mu,\mathrm{raw}}\var_{\C}$ for the ring with good finite
$\widehat\mu$-actions and the equivariant scissor relations.
Here \emph{good} means that the action factors through some $\mu_N$
and every orbit is contained in an affine open subset.
Write $\gmv$ for its quotient by the ideal generated by
\[
 [X\times(\A^r,\rho)]-[X\times\A^r],
\]
for every good equivariant $X$ and every finite linear action
$\rho$; the action on the second affine factor on the right is
trivial. Goodness and this additional relation are separate
conditions.

There is a canonical isomorphism
\[
 \Theta:\sggk\RES\xrightarrow{\ \sim\ }\gmv .
\]
We construct it at a finite Kummer level and prove both
directions of the comparison below.
\end{nota}

\begin{proof}
\emph{Finite-level coordinates and the raw comparison.}
For $q\in\Q$, let $L_q=q^\sharp\sqcup\{0_q\}$ be the
one-dimensional residue vector space whose nonzero part is
$q^\sharp=\vrv^{-1}(q)$. The notation $0_q$ records its zero
in that fixed vector space, encoded using the existing RV
zero/infinity element. It introduces no second RV zero.

For $U\in\RES$, split its finite value support and its zero
coordinate positions. Each piece can be regarded as a definable
subset of $\prod_{j=1}^r L_{a_j/N}$ for a common $N>0$ and
integers $a_j$. Choose one root $\tau_N^N=t$ and put
$r_N=\rv(\tau_N)$. The coordinate map
\[
 g_{r_N}(u_1,\ldots,u_r)
   =(u_1/r_N^{a_1},\ldots,u_r/r_N^{a_r})
\]
takes this product to $\C^r$, with $0_q$ mapped to $0$.
Residue-field purity makes $X_U=g_{r_N}(U)$ constructible.
The action of $\mu_N$ on its coordinates is
\[
 (x_j)_j\longmapsto(\zeta^{a_j}x_j)_j .
\]
The set $X_U$ is invariant because $U$ is defined over $F$.
Replacing $r_N$ by $\eta r_N$ multiplies the $j$th coordinate
by $\eta^{-a_j}$, an equivariant algebraic isomorphism.
Enlarging $N$ only inflates the action along
$\mu_M\to\mu_N$ and changes coordinates by such constants.
Thus the resulting equivariant constructible class is independent
of both choices. Products use a common Kummer level, so the
assignment respects products as well as scissors.

We now prove that every equivariant constructible set and
bijection is obtained in this way. A good finite action admits a
finite invariant affine stratification: intersect the translates
of an affine neighborhood of each orbit, take a finite such
cover, and successively remove the preceding opens. These
intersections are affine because the varieties are separated.
For an affine $\mu_N$-variety $X$, take algebra generators of
$\C[X]$, span their finite orbits, and diagonalize that finite
dimensional cyclic representation. Its eigenvectors still
generate the algebra. They give an equivariant closed embedding
\[
 X\hookrightarrow\A^r,\qquad
 \zeta\cdot(x_j)_j=(\zeta^{a_j}x_j)_j
\]
for suitable integers $a_j$.

Let
\[
 T_N=\{r\in\RV:r^N=\rv(t)\}.
\]
For any invariant constructible
$X\subseteq\A^r$ with these coordinate weights, define
\[
 U_X=\left\{u\in\prod_jL_{a_j/N}:
       \exists r\in T_N\quad g_r(u)\in X\right\}.
\]
This is an original-language $F$-definable RES set.
Every two elements of $T_N$ differ by an element of $\mu_N$.
Invariance of $X$ therefore implies
$g_r(U_X)=X$ for every $r\in T_N$.
This proves surjectivity without adding a root to the base.

The same construction applies to the graph of an equivariant
constructible bijection $f:X\to Y$, using one common level:
\[
 G_f=\{(u,v):\exists r\in T_N,\
              (g_r(u),g_r(v))\in\operatorname{Graph}(f)\}.
\]
Equivariance makes the condition independent of $r$, so $G_f$
is the graph of an $F$-definable bijection.
Conversely, framing the graph of an $F$-definable RES bijection
gives an equivariant constructible bijection.
In characteristic zero such a bijection is piecewise an
algebraic isomorphism: its generic graph over each irreducible
piece has degree one, hence a rational inverse on a dense open;
remove those opens and induct on dimension. Intersecting the
finitely many translates makes the partition invariant.
Consequently these graph constructions also respect the usual
equivariant scissor presentation. They show independence of the
chosen affine embeddings and stratifications and give inverse
maps on classes. We have proved an isomorphism
\[
 \Theta_{\mathrm{raw}}:\ggk\RES\xrightarrow{\ \sim\ }\ggk^{\widehat\mu,\mathrm{raw}}\var_{\C} .
\]

\emph{The quotient.}
For $q=a/N$, the space $L_q$ is sent to the affine line with
character $\rho_a(\zeta)x=\zeta^a x$. Its zero is a fixed point.
Hence
\[
 \Theta_{\mathrm{raw}}
 \bigl([q^\sharp]-[\G_m]\bigr)
 =
 [\A^1,\rho_a]-[\A^1,\mathrm{triv}].
\]
Every finite character of $\widehat\mu$ has this form.
Every finite linear cyclic action is a sum of characters, so its
class is a product of these line classes. By expanding the
difference of two such products, all the linear-action relations,
also after multiplication by an arbitrary $[X]$, are generated
by the displayed line differences. The two quotient ideals
therefore correspond exactly. This proves the asserted
isomorphism $\Theta$.

For later computations, let
$h_{M_i}(U_i)=V_i\times I_i^\sharp$ be a monomial decomposition
as in Lemma~\ref{coh:decom}, with $\ell_i$ free RV coordinates.
The original-language construction of $\bb E_b$ in Section~\ref{sec:retract} gives
\begin{equation}\label{twist:C}
 (\Theta\circ\bb E_b)([\bm U])
   =\sum_i\chi_b(I_i)\,\Theta([V_i])[\G_m]^{\ell_i}.
\end{equation}
The finite-to-one label and its degree are retained as in the
product construction of Section~\ref{sec:setup}; the displayed target forgets
that grading. The formula is an evaluation of the already defined
retraction, so it is independent of the decomposition.
\end{proof}

\begin{nota}\label{mov:vol}
The motivic volume operator is
\[
 \Vol=\Theta\circ\bb E_b\circ\int:
 \ggk\VF_*\longrightarrow\gmv .
\]
Here $\bb E_b$ is applied after the quotient by $(\bm P-1)$,
on which it is already defined by Section~\ref{sec:retract}.
\end{nota}

\begin{cor}\label{cor:global:vanishing}
Let $F=\C((t))$ and let $f:A\to B$ be a morphism of
varieties over $\C$.
Interpret these varieties in a sufficiently saturated
ambient model $\mathbb U$ of $\mathrm{ACVF}(0,0)$, using
finite affine charts.
For $b\in B(\mathbb U)$ put $A_b=f^{-1}(b)$ and
$\mathbb S_b=F\langle b\rangle$, with $b$ represented by its
VF chart coordinates.
If, for every $b\in B(\mathbb U)$,
\[
 \mathbb E_{\mathrm b}^{\mathbb S_b}\!
       \left(\int[A_b]\right)=0
       \quad\text{in }R(\mathbb S_b),
\]
then $\Vol([A])=0$ in $\gmv$.
\end{cor}

\begin{proof}
Choose finite disjoint constructible pieces subordinate to
affine covers of $A$ and $B$ and encode their finite tags
by distinct elements of $\C$.
This realizes $f$ as a definable map between VF tuples.
Changing charts gives definably equivalent parameter
bases and transports the fiber equality.

Both $F$ and each $\mathbb S_b$ in these coordinates are
VF-generated. Thus the integrals of Section~\ref{sec:setup} are
available over these bases.
The lifting theorem supplies positive RV representatives
of $A_b$ with actual $\mathbb S_b$-definable lift bijections.
Apply Theorem~\ref{fubini} to $A\to B$ and the empty family
$\varnothing\to B$.
The displayed hypothesis is precisely the required
equality of their fiber retractions, so
\[
 \mathbb E_{\mathrm b}^{F}\!
       \left(\int[A]\right)=0
       \quad\text{in }R(F).
\]
Applying the global comparison $\Theta$ gives
$\Vol([A])=0$.
\end{proof}
